\documentclass[11pt, a4paper]{article}

% bfw-style.tex — gemeinsame Style-Definitionen für alle Handouts/Aufgabenhefte
% Voraussetzung: Kompilation mit xelatex (wegen fontspec)

\usepackage[ngerman]{babel}
\usepackage{geometry}
\geometry{a4paper, top=2.2cm, bottom=2.2cm, left=2.3cm, right=2.3cm}

\usepackage{fontspec}
% Fallback-Kette: Source Sans Pro -> Calibri -> Segoe UI -> Latin Modern Sans
% (Latin Modern ist immer mit MiKTeX vorhanden)
\IfFontExistsTF{Source Sans Pro}{%
  \setmainfont{Source Sans Pro}[Ligatures=TeX]%
}{\IfFontExistsTF{Calibri}{%
  \setmainfont{Calibri}[Ligatures=TeX]%
}{\IfFontExistsTF{Segoe UI}{%
  \setmainfont{Segoe UI}[Ligatures=TeX]%
}{%
  \setmainfont{Latin Modern Sans}[Ligatures=TeX]%
}}}
\IfFontExistsTF{Source Code Pro}{%
  \setmonofont{Source Code Pro}[Scale=0.92]%
}{\IfFontExistsTF{Consolas}{%
  \setmonofont{Consolas}[Scale=0.92]%
}{%
  \setmonofont{Latin Modern Mono}[Scale=0.92]%
}}

\usepackage{xcolor}
% Farbpalette: hell, freundlich, modern
\definecolor{bfwPrimary}{HTML}{0EA5A4}    % Petrol/Türkis - Primär
\definecolor{bfwPrimaryDark}{HTML}{0F766E} % dunkler Petrol für Headings
\definecolor{bfwAccent}{HTML}{F59E0B}     % Amber - Akzent
\definecolor{bfwSuccess}{HTML}{10B981}    % Grün - Erfolg / Lösung
\definecolor{bfwWarn}{HTML}{F97316}       % Orange - Achtung
\definecolor{bfwInk}{HTML}{0F172A}        % fast Schwarz
\definecolor{bfwInkSoft}{HTML}{475569}    % gedämpftes Grau
\definecolor{bfwBgSoft}{HTML}{F8FAFC}     % off-white
\definecolor{bfwBgTeal}{HTML}{ECFEFF}     % helles Türkis
\definecolor{bfwBgAmber}{HTML}{FEF3C7}    % helles Gelb
\definecolor{bfwBgRose}{HTML}{FFE4E6}     % helles Rosé für Eselsbrücken

\usepackage[colorlinks=true, linkcolor=bfwPrimaryDark, urlcolor=bfwPrimaryDark]{hyperref}
\usepackage{enumitem}
\setlist[itemize]{leftmargin=*, itemsep=2pt, topsep=2pt}
\setlist[enumerate]{leftmargin=*, itemsep=2pt, topsep=2pt}

\usepackage[most]{tcolorbox}
\usepackage{booktabs}
\usepackage{tikz}
\usetikzlibrary{decorations.pathmorphing, positioning, shapes.geometric, arrows.meta}

% Merkkästchen — wichtige Definition / Kernpunkt
\newtcolorbox{merksatz}[1][]{%
  enhanced, breakable,
  colback=bfwBgTeal,
  colframe=bfwPrimary,
  boxrule=0.6pt,
  arc=4pt,
  left=10pt, right=10pt, top=8pt, bottom=8pt,
  fonttitle=\bfseries\color{bfwPrimaryDark},
  title={\faLightbulb\ Merksatz},
  #1
}

% Eselsbrücke — Mnemoniks
\newtcolorbox{eselsbruecke}[1][]{%
  enhanced, breakable,
  colback=bfwBgRose,
  colframe=bfwAccent,
  boxrule=0.6pt,
  arc=4pt,
  left=10pt, right=10pt, top=8pt, bottom=8pt,
  fonttitle=\bfseries\color{bfwWarn},
  title={Eselsbrücke},
  #1
}

% Beispiel-Box
\newtcolorbox{beispiel}[1][]{%
  enhanced, breakable,
  colback=bfwBgSoft,
  colframe=bfwInkSoft,
  boxrule=0.4pt,
  arc=3pt,
  left=10pt, right=10pt, top=8pt, bottom=8pt,
  fonttitle=\bfseries\color{bfwInk},
  title={Beispiel},
  #1
}

% Lösungs-Box (für Aufgabenhefte)
\newtcolorbox{loesung}[1][]{%
  enhanced, breakable,
  colback=bfwBgAmber!40,
  colframe=bfwSuccess,
  boxrule=0.6pt,
  arc=3pt,
  left=10pt, right=10pt, top=8pt, bottom=8pt,
  fonttitle=\bfseries\color{bfwSuccess!70!black},
  title={Lösung},
  #1
}

% Glossar-Eintrag
\newcommand{\glossareintrag}[2]{%
  \par\noindent\textbf{\color{bfwPrimaryDark}#1} \enspace #2\par\smallskip
}

% Section-Stil — etwas farbiger als Standard
\usepackage{titlesec}
\titleformat{\section}
  {\Large\bfseries\color{bfwPrimaryDark}}
  {\thesection}{0.6em}{}
\titleformat{\subsection}
  {\large\bfseries\color{bfwPrimary}}
  {\thesubsection}{0.5em}{}
\titleformat{\subsubsection}
  {\normalsize\bfseries\color{bfwInk}}
  {\thesubsubsection}{0.4em}{}

% TikZ-Stil "skizzenhaft" — etwas weicher als Rough.js, aber stabil
\tikzset{
  roughbox/.style = {
    draw=bfwPrimaryDark, thick, fill=bfwBgTeal,
    rounded corners=4pt, inner sep=6pt
  },
  roughaccent/.style = {
    draw=bfwAccent, thick, fill=bfwBgAmber,
    rounded corners=4pt, inner sep=6pt
  },
  roughline/.style = {
    draw=bfwInkSoft, thick, -{Stealth[length=2.5mm]}
  }
}

% Bullet-Symbole (bewusst kein FontAwesome — vermeidet Font-Installations-Probleme)
\newcommand{\faLightbulb}{$\bullet$}
\newcommand{\faBookmark}{$\bullet$}

% Header/Footer
\usepackage{fancyhdr}
\pagestyle{fancy}
\fancyhf{}
\renewcommand{\headrulewidth}{0pt}
\fancyfoot[L]{\small\color{bfwInkSoft}\thepage}
\fancyfoot[R]{\small\color{bfwInkSoft} BFW M-064 Beschaffung im Büromanagement}


\title{\vspace*{-1.5cm}\Huge\bfseries\color{bfwPrimaryDark}Aufgabenheft Tag~4\\[0.4em]
  \large\color{bfwInkSoft}\textnormal{Arbeitsschutz \& Gesundheit am Arbeitsplatz --- Übungen im IHK-Klausurformat}}
\author{\textsf{Büroprozesse gestalten und Arbeitsvorgänge organisieren \textperiodcentered{} M-067}\\
\textsf{\small\copyright\ Can Siebert\,\textperiodcentered\,2026}}
\date{Selbstlernmaterial \textperiodcentered{} 02.07.2026}

\begin{document}
\maketitle
\thispagestyle{empty}

\begin{merksatz}[title={\bfseries So nutzen Sie dieses Aufgabenheft}]
Bearbeiten Sie alle Aufgaben zunächst \emph{ohne} in die Lösungen zu schauen. Schreiben
Sie Ihre Antworten in vollständigen Sätzen --- die IHK-Prüfung verlangt formulierte
Begründungen, nicht bloße Stichworte. Beachten Sie die \emph{Operatoren} (nennen,
erläutern, beurteilen, begründen) --- sie geben den geforderten Antwort-Tiefegrad vor.
Erst danach öffnen Sie den Lösungsteil ab Seite~\pageref{loesungen} und vergleichen.
\end{merksatz}

\tableofcontents
\vspace{1em}
\hrule

\section{Aufgaben}

\subsection*{Aufgabe~1 --- Rechtsgrundlagen des Arbeitsschutzes (10 Punkte)}

\noindent\textbf{\color{bfwAccent}YouTube-Vorbereitung:}\enspace \href{https://www.youtube.com/results?search_query=Arbeitsschutzgesetz+einfach+erklaert}{Arbeitsschutzgesetz einfach erklärt} (\url{https://www.youtube.com/results?search_query=Arbeitsschutzgesetz+einfach+erklaert})

\medskip
Sie sind in der Verwaltung der \emph{Duisdorfer BüroKonzept KG} tätig. Die Geschäftsführung
bittet Sie um einen Überblick über die wichtigsten Arbeitsschutzgesetze.

\begin{enumerate}[label=(\alph*)]
  \item \textbf{Erläutern} Sie das Ziel des Arbeitsschutzgesetzes (ArbSchG). \hfill (3~P.)
  \item \textbf{Nennen} Sie zwei Regelungsbereiche des Arbeitszeitgesetzes (ArbZG). \hfill (2~P.)
  \item \textbf{Beschreiben} Sie, wozu das Arbeitssicherheitsgesetz (ASiG) den Arbeitgeber verpflichtet. \hfill (3~P.)
  \item \textbf{Beurteilen} Sie die Aussage: \emph{„Für die Gesundheit am Arbeitsplatz ist allein der Arbeitgeber verantwortlich.”} \hfill (2~P.)
\end{enumerate}

\subsection*{Aufgabe~2 --- Verantwortung und Aufsicht (12 Punkte)}

\noindent\textbf{\color{bfwAccent}YouTube-Vorbereitung:}\enspace \href{https://www.youtube.com/results?search_query=Berufsgenossenschaft+Aufgaben}{Berufsgenossenschaft --- Aufgaben} (\url{https://www.youtube.com/results?search_query=Berufsgenossenschaft+Aufgaben})

\medskip
\begin{enumerate}[label=(\alph*)]
  \item \textbf{Nennen} Sie zwei Pflichten des Arbeitgebers und zwei Pflichten des Arbeitnehmers im Arbeitsschutz. \hfill (4~P.)
  \item \textbf{Erläutern} Sie drei Aufgaben der Berufsgenossenschaften. \hfill (3~P.)
  \item \textbf{Beschreiben} Sie die Aufgaben und Befugnisse des Gewerbeaufsichtsamtes. \hfill (3~P.)
  \item \textbf{Erklären} Sie den Unterschied zwischen Betriebsarzt und Sicherheitsbeauftragtem. \hfill (2~P.)
\end{enumerate}

\subsection*{Aufgabe~3 --- Gefährdungsbeurteilung und Unterweisung (12 Punkte)}

\noindent\textbf{\color{bfwAccent}YouTube-Vorbereitung:}\enspace \href{https://www.youtube.com/results?search_query=Gefaehrdungsbeurteilung+erklaert}{Gefährdungsbeurteilung erklärt} (\url{https://www.youtube.com/results?search_query=Gefaehrdungsbeurteilung+erklaert})

\medskip
In der \emph{Duisdorfer BüroKonzept KG} wird ein neuer Bildschirmarbeitsplatz eingerichtet.

\begin{enumerate}[label=(\alph*)]
  \item \textbf{Erläutern} Sie, wozu die Gefährdungsbeurteilung nach \S~5 ArbSchG dient. \hfill (3~P.)
  \item \textbf{Nennen} Sie drei Zeitpunkte, zu denen eine Gefährdungsbeurteilung durchzuführen ist. \hfill (3~P.)
  \item \textbf{Erläutern} Sie den Unterschied zwischen Unterweisung und Betriebsanweisung. \hfill (3~P.)
  \item \textbf{Beurteilen} Sie, warum psychische Belastungen heute Teil der Gefährdungsbeurteilung sind. \hfill (3~P.)
\end{enumerate}

\subsection*{Aufgabe~4 --- Belastung, Beanspruchung und Stress (15 Punkte)}

\noindent\textbf{\color{bfwAccent}YouTube-Vorbereitung:}\enspace \href{https://www.youtube.com/results?search_query=Eustress+Distress+Unterschied}{Eustress und Distress} (\url{https://www.youtube.com/results?search_query=Eustress+Distress+Unterschied})

\medskip
Herr Unger aus der Personalabteilung klagt über ständige Kopfschmerzen. Er hat viel Arbeit,
wird ständig unterbrochen und macht selten Pausen.

\begin{enumerate}[label=(\alph*)]
  \item \textbf{Erläutern} Sie den Unterschied zwischen Belastung und Beanspruchung. \hfill (4~P.)
  \item \textbf{Grenzen} Sie Eustress von Distress \textbf{ab}. \hfill (4~P.)
  \item \textbf{Nennen} Sie vier Symptome von negativem Stress. \hfill (4~P.)
  \item \textbf{Nennen} Sie drei Maßnahmen, mit denen Herr Unger seinen Stress abbauen kann. \hfill (3~P.)
\end{enumerate}

\subsection*{Aufgabe~5 --- Burnout und Work-Life-Balance (12 Punkte)}

\noindent\textbf{\color{bfwAccent}YouTube-Vorbereitung:}\enspace \href{https://www.youtube.com/results?search_query=Burnout+Symptome+Ursachen}{Burnout --- Symptome und Ursachen} (\url{https://www.youtube.com/results?search_query=Burnout+Symptome+Ursachen})

\medskip
\begin{enumerate}[label=(\alph*)]
  \item \textbf{Erläutern} Sie, was unter Burnout zu verstehen ist. \hfill (3~P.)
  \item \textbf{Nennen} Sie drei typische Burnout-Symptome. \hfill (3~P.)
  \item \textbf{Beschreiben} Sie, welche Personen besonders gefährdet sind. \hfill (3~P.)
  \item \textbf{Erläutern} Sie zwei Maßnahmen, mit denen Unternehmen die Work-Life-Balance fördern. \hfill (3~P.)
\end{enumerate}

\subsection*{Aufgabe~6 --- Mobbing und betriebliches Gesundheitsmanagement (10 Punkte)}

\noindent\textbf{\color{bfwAccent}YouTube-Vorbereitung:}\enspace \href{https://www.youtube.com/results?search_query=Mobbing+am+Arbeitsplatz+Massnahmen}{Mobbing am Arbeitsplatz} (\url{https://www.youtube.com/results?search_query=Mobbing+am+Arbeitsplatz+Massnahmen})

\medskip
\begin{enumerate}[label=(\alph*)]
  \item \textbf{Nennen} Sie zwei Kriterien, an denen Mobbing von einem gewöhnlichen Konflikt zu unterscheiden ist. \hfill (2~P.)
  \item \textbf{Ordnen} Sie für jeden Fall die Ursachenebene zu (Täter, Opfer, Rahmenbedingungen) \textbf{und begründen} Sie: \hfill (3~P.)
    \begin{enumerate}[label=\arabic*., leftmargin=*]
      \item Ein Vorgesetzter will durch ständiges Kritisieren seine eigene Kompetenz demonstrieren.
      \item Hoher Leistungsdruck und umständliche Dienstwege erschweren das Arbeiten.
      \item Ein neuer Mitarbeiter kennt die Abläufe noch nicht und tritt zu forsch auf.
    \end{enumerate}
  \item \textbf{Nennen} Sie drei Maßnahmen, die ein Mobbing-Opfer ergreifen kann. \hfill (3~P.)
  \item \textbf{Beurteilen} Sie, warum Suchtprävention in das betriebliche Gesundheitsmanagement gehört. \hfill (2~P.)
\end{enumerate}

\vspace{1em}
\noindent\textbf{Punkte gesamt: 71}

\newpage

\section{Lösungen}\label{loesungen}

\subsection*{Lösung Aufgabe~1}
\begin{loesung}
\textbf{(a)} Das Arbeitsschutzgesetz (ArbSchG) enthält Maßnahmen, die die Sicherheit und
den Gesundheitsschutz der Beschäftigten bei der Arbeit sichern und verbessern sollen. Es
geht um Gefahren, die \emph{während} und \emph{durch} die Arbeit entstehen; \S~4 nennt
allgemeine Grundsätze (z.\,B.\ Gefahren an der Quelle bekämpfen).

\textbf{(b)} Das ArbZG regelt z.\,B.\ die tägliche Höchstarbeitszeit, die
Mindestruhepausen, den besonderen Schutz von Nachtarbeitern und das grundsätzliche
Arbeitsverbot an Sonn- und Feiertagen.

\textbf{(c)} Das ASiG verpflichtet den Arbeitgeber, Betriebsärzte und Fachkräfte für
Arbeitssicherheit zu bestellen. Diese unterstützen ihn beim Arbeitsschutz und bei der
Unfallverhütung.

\textbf{(d)} Die Aussage ist falsch. Die Verantwortung liegt \emph{sowohl} beim Arbeitgeber
\emph{als auch} beim Arbeitnehmer. Der Arbeitgeber schafft sichere Bedingungen; der
Arbeitnehmer muss sich so verhalten, dass er Gefährdungen vorbeugt, Gefahren meidet und auf
Schwachstellen hinweist.
\end{loesung}

\subsection*{Lösung Aufgabe~2}
\begin{loesung}
\textbf{(a)} Arbeitgeber: für geeignete Arbeitsbedingungen und Arbeitsmittel sorgen, die
Einhaltung von Arbeitsschutz- und Unfallverhütungsvorschriften gewährleisten. Arbeitnehmer:
sich so verhalten, dass Gesundheitsgefährdungen vorgebeugt wird; Gefahren meiden und den
Arbeitgeber auf Schwachstellen hinweisen.

\textbf{(b)} Berufsgenossenschaften überwachen die Verhütung von Arbeitsunfällen,
Berufskrankheiten und arbeitsbedingten Gesundheitsgefahren, kontrollieren die Erste Hilfe,
beraten Unternehmen und Beschäftigte und erlassen die Berufsgenossenschaftlichen
Vorschriften (BGV).

\textbf{(c)} Das Gewerbeaufsichtsamt kontrolliert die Einhaltung der Vorschriften zum
Arbeits-, Umwelt- und Verbraucherschutz, darf Arbeitsstätten jederzeit betreten und
besichtigen, kann Anordnungen erteilen, Betriebe stilllegen und Bußgelder verhängen
(\S~35 GewO).

\textbf{(d)} Der Betriebsarzt ist eine arbeitsmedizinische Fachkraft, die berät, untersucht
und die Durchführung des Arbeitsschutzes beobachtet. Der Sicherheitsbeauftragte ist eine
vom Unternehmen schriftlich bestellte Person (SGB), die den Arbeitsschutz vor Ort als
Bindeglied unterstützt.
\end{loesung}

\subsection*{Lösung Aufgabe~3}
\begin{loesung}
\textbf{(a)} Der Arbeitgeber muss durch die Beurteilung der mit der Arbeit verbundenen
Gefährdung ermitteln, welche Maßnahmen des Arbeitsschutzes erforderlich sind (\S~5 ArbSchG).
Berücksichtigt werden u.\,a.\ Gestaltung der Arbeitsstätte, Einwirkungen, Arbeitsmittel,
Arbeitsverfahren und Arbeitszeit, Qualifikation sowie die psychische Belastung.

\textbf{(b)} Vor Aufnahme der Tätigkeiten, beim Einrichten und Betreiben von Arbeitsstätten
und vor der erstmaligen Verwendung eines Arbeitsmittels. Sie ist zu aktualisieren, wenn
neue Arbeitsplätze oder Verfahren hinzukommen.

\textbf{(c)} Die Unterweisung (\S~12 ArbSchG) ist eine mündliche, arbeitsplatzbezogene
und regelmäßige Belehrung der Beschäftigten über Gefahren und richtiges Verhalten; sie wird
schriftlich dokumentiert und unterschrieben. Die Betriebsanweisung ist ein schriftliches,
ausgehängtes Dokument zu Gefahren, Schutzmaßnahmen und Verhaltensregeln.

\textbf{(d)} Seit Oktober 2013 werden psychische Belastungen ausdrücklich als
Gesundheitsgefährdung angesehen. Da Stress, Termindruck oder Konflikte krank machen können,
muss der Arbeitgeber sie ebenso wie körperliche Gefahren feststellen und Maßnahmen ergreifen.
\end{loesung}

\subsection*{Lösung Aufgabe~4}
\begin{loesung}
\textbf{(a)} Belastung ist ein neutraler Begriff für alle Einflüsse, die von außen auf den
Menschen einwirken. Beanspruchung ist die individuelle Wirkung dieser Belastung im Menschen
und hängt von Faktoren wie Belastungsstärke, Alter, Konstitution, Qualifikation,
Handlungsspielraum und Stressumgang ab.

\textbf{(b)} Eustress (positiver Stress) erhöht die Aufmerksamkeit, fördert die
Leistungsfähigkeit ohne zu schaden, motiviert und steigert die Produktivität. Distress
(negativer Stress) entsteht, wenn Stress zu häufig oder ohne körperlichen Ausgleich
auftritt; er überfordert und ist bedrohlich.

\textbf{(c)} Vier Symptome (Beispiele): verminderte Leistungs- und Konzentrationsfähigkeit,
Schlafstörungen, Magen-Darm-Probleme, Gereiztheit, Bluthochdruck.

\textbf{(d)} Maßnahmen: Sport treiben, Freizeitstress vermeiden, Erholungsphasen einplanen
(bevor man sich gestresst fühlt), ausreichend schlafen, Entspannungstechniken (Yoga,
Meditation); zusätzlich Prioritäten setzen und Störfaktoren minimieren.
\end{loesung}

\subsection*{Lösung Aufgabe~5}
\begin{loesung}
\textbf{(a)} Burnout („ausbrennen”) ist ein Zustand völliger physischer und psychischer
Erschöpfung aufgrund meist beruflicher Überlastung. Es gilt als Vorstufe zur Depression,
ist keine eigenständige Krankheit mit klarer Diagnose und entsteht schleichend.

\textbf{(b)} Drei Symptome (Beispiele): chronische Müdigkeit, Konzentrationsprobleme,
Schlafstörungen, Erschöpfung, Kopf- oder Rückenschmerzen, Verspannungen.

\textbf{(c)} Besonders gefährdet sind Menschen mit hohen Leistungsansprüchen,
Perfektionisten und Personen, die sich vor allem über ihren Job definieren; Ursachen sind
oft großes Engagement und übertriebener Ehrgeiz.

\textbf{(d)} Maßnahmen (Beispiele): flexible Arbeitszeitmodelle (Teilzeit, Gleitzeit,
Vertrauensarbeitszeit), veränderte Arbeitsorganisation (Job Rotation, Job Enlargement),
betriebliches Gesundheitsmanagement, Unterstützung bei Alltagsproblemen, Weiterbildung im
Stress- und Zeitmanagement.
\end{loesung}

\subsection*{Lösung Aufgabe~6}
\begin{loesung}
\textbf{(a)} Zwei Kriterien (Beispiele): Der Konflikt hat sich verfestigt; die Würde des
Menschen wird angegriffen; die Parteien begegnen sich nicht auf Augenhöhe; der Unterlegene
kommt nicht aus eigener Kraft heraus. Zudem sind die Handlungen wiederholt und mit
Schädigungsabsicht.

\textbf{(b)} Zuordnung mit Begründung:
\begin{enumerate}[label=\arabic*., itemsep=0pt]
  \item \textbf{Täter} --- der Mobber verfolgt ein Ziel (Kompetenz demonstrieren) bzw.\ handelt aus Angst.
  \item \textbf{Rahmenbedingungen} --- Leistungsdruck und umständliche Dienstwege sind Umstände der Arbeitswelt, die Mobbing begünstigen.
  \item \textbf{Opfer} --- ein neuer Mitarbeiter ist besonders gefährdet und liefert (unabsichtlich) einen Anlass.
\end{enumerate}

\textbf{(c)} Drei Maßnahmen (Beispiele): selbstkritische Prüfung, den Täter ansprechen und
Grenzen setzen, soziales Netzwerk/Selbstbewusstsein stärken, Mobbing-Tagebuch führen,
Rechtsweg beschreiten oder Arbeitsplatz wechseln. Wichtig: früh reagieren.

\textbf{(d)} Arbeitsbedingungen wie hoher Leistungsdruck, ungünstige Arbeitszeiten und
schwieriges Betriebsklima erhöhen das Suchtrisiko. Suchtmittel gefährden die
Arbeitssicherheit (Konzentrationsstörungen, längere Reaktionszeit, Unfälle). Deshalb
gehören Information, Schulung und Beratung als Suchtprävention zum BGM.
\end{loesung}

\end{document}
